% Generated draft: requires mathematical and source review.
\documentclass[11pt]{article}
\usepackage[margin=25mm]{geometry}
\usepackage{fontspec}
\setmainfont{Latin Modern Roman}
\usepackage{amsmath,amssymb,mathtools,unicode-math}
\setmathfont{Latin Modern Math}
\usepackage{enumitem,needspace,xurl,hyperref}
\setlength{\parindent}{0pt}
\setlength{\parskip}{5pt}
\setlength{\emergencystretch}{3em}
\newcommand{\N}{\mathbb{N}}
\newcommand{\Z}{\mathbb{Z}}
\newcommand{\Q}{\mathbb{Q}}
\newcommand{\R}{\mathbb{R}}
\newcommand{\C}{\mathbb{C}}
\newcommand{\F}{\mathbb{F}}
\newcommand{\A}{\mathbb{A}}
\newcommand{\PP}{\mathbb{P}}
\newcommand{\E}{\mathbb{E}}
\newcommand{\eps}{\varepsilon}
\newcommand{\dd}{\,\mathrm{d}}
\begin{document}
\section*{UnsolvedMath ID: 700948}
\textbf{Source problem code:} catalogue entry
\subsection{Definitions and mathematical statement}
For a group \(G\), let \(\operatorname{Sym}(G)\) be the group of all permutations of its underlying set, and let \(\lambda:G\to\operatorname{Sym}(G)\) be its left regular representation, defined by \(\lambda(g)(x)=gx\). A permutation group \(H\leq\operatorname{Sym}(G)\) is primitive if it is transitive and the only \(H\)-invariant equivalence relations on \(G\) are equality and the universal relation. For \(|G|\geq2\), it is doubly transitive if, for all \(x,y,u,v\in G\) with \(x\ne y\) and \(u\ne v\), there exists \(h\in H\) such that \(h(x)=u\) and \(h(y)=v\).

For \(|G|\geq2\), the supplied definition of a \(B\)-group is therefore
\[
\forall H\leq\operatorname{Sym}(G),\qquad
\bigl(\lambda(G)\leq H\ \text{and }H\text{ is primitive}\bigr)
\Longrightarrow H\text{ is doubly transitive}.
\]
The source asks: does there exist a countable \(B\)-group?

The meaning of ``countable'' is unresolved in the supplied record. If finite groups are included, the cardinality condition is \(|G|\leq\aleph_0\); if countably infinite groups are intended, it is \(|G|=\aleph_0\). These are distinct formulations, and the source does not select between them. The source also gives no convention for primitivity or double transitivity in degree one; the explicit definitions above apply to \(|G|\geq2\), leaving the treatment of the trivial group unresolved.
\subsection{Short English statement}
P. J. Cameron asks whether a countable group can have the property that every primitive permutation group containing its regular representation is doubly transitive. The source does not clarify whether ``countable'' includes finite groups or how the trivial group is treated.
\subsection{Sources}
\begin{itemize}
\item Local UnsolvedMath registry, permanent ID \(700948\), exact source code \texttt{catalogue entry}, title \emph{Kourovka Notebook Problem 13.30}. The supplied registry labels the problem open; this label is not an independently verified mathematical or literature assessment.
\item P. J. Cameron, problem 13.30 in \emph{The Kourovka Notebook (Unsolved Problems in Group Theory)}, 13th issue (1995), as attributed in the supplied record.
\item The supplied background also identifies the 21st issue edition (2026), arXiv:1401.0300v48. The collection description names E. I. Khukhro and V. D. Mazurov as editors. This reference information is supplied attribution and has not been independently checked.
\end{itemize}
\end{document}
